% Study sheet: Volume I, Chapter 4, Data Engineering.
% Numbers: \val{...} from data_engineering_studysheet_numbers.tex, written by data_engineering_studysheet_numbers.py.
\documentclass{studysheet}
\input{data_engineering_studysheet_numbers}
\sheetvolume{1}
\sheetchapter{4}{Data Engineering}
\sheetsources{Numbers from the chapter's keyword-spotting and ResNet-50 scenarios in mlsysim; assumed inputs are marked.}

\begin{document}
% ===================================================================== FRONT: LEARN
\frontside
\begin{bigidea}
A model can learn only what its data contains, and it trains only as fast as its data arrives.
Data is therefore \hl{code}, to be versioned, tested and released, and \hl{freight}, to be stored,
delivered as fast as the accelerator consumes it.
\end{bigidea}

\begin{leftcol}
\sheetsec{Key ideas}
\begin{ideas}
\idea{Data is code.}{Changing the dataset changes the model with no code change. Splits are trust
  boundaries: put each speaker or household wholly in one split, or the test score measures memory.}
\idea{Defects cascade silently.}{A wrong field passes schemas and trains fine while every later stage
  builds on it. Check at ingestion: rules on each record, statistics on each batch.}
\idea{Labels are the specification.}{A wrong label redefines ``correct''; labels are often the largest
  data cost. Check agreement beyond chance and known answers.}
\idea{Fit once, replay everywhere.}{Fit scaling on training data only, store it with the model, use
  only inputs serving has. Else leakage or silent \emph{training-serving skew}.}
\idea{Training runs no faster than its data.}{Storage, host decoding and bytes per sample each cap the
  supply. Large datasets pull computation to them (\emph{data gravity}).}
\idea{Drift warns; labels measure.}{Drift comes from the world, skew from the stack. Logs show what
  the model caught, not what it missed.}
\end{ideas}

\vspace{11pt}
\sheetsec{Which failure is it?}
\begin{sheettable}{@{}l L L@{}}
\thead{Cause} & \thead{What changed} & \thead{Evidence that separates it}\\ \midrule
Skew & the stack: features computed differently & replay fixed raw inputs through both paths\\
Covariate drift & the world: inputs moved & input statistics vs.\ training (PSI, K-S); old vs.\ new cohorts on one path and rubric\\
Rubric change & the yardstick: grading rules & rescore fixed predictions under both rubrics\\
Concept drift & the input-to-label link & fresh labels only\\
Schema change & an upstream field & contract check at ingestion\\
\end{sheettable}
\vspace{3pt}
{\fontsize{8.6}{11}\selectfont\color{muted}Population stability index (PSI) below \val{psiwarn}: stable;
above \val{psicrit}: investigate.\par}
\end{leftcol}%
\begin{rightcol}
\sheetsec{Formulas you will use}
\fcard{Data supply bound}%
  {\begin{gathered}\TrainingSampleThroughput \le \min(\ComputeLimitedSampleRate,\,\DataSupplyRate)\\[2pt]
   \DataSupplyRate = \frac{\StorageBandwidth(1-\OperationalOverheadFraction)}{\SampleStorageSize}\end{gathered}}%
  {Samples trained per second: the smaller of what compute uses and what storage delivers (bandwidth
   $\StorageBandwidth$, overhead share $\OperationalOverheadFraction$, bytes per sample $\SampleStorageSize$).
   Host decoding also caps supply; use the lower rate. Feeding loss $= 1 - $ supply $\div$ demand
   when supply is short; $\OperationalOverheadFraction = 0$ unless given.}%
  {Storage-bound: faster tier, sequential shards, denser codec, fewer bytes per sample.
   Decode-bound: more workers, hardware decode, a faster-to-decode codec. Storing decoded samples
   skips decoding but multiplies bytes and fixes the stored crops.}{}
\fcard{False-wake budget}%
  {\text{false-positive rate per window} \le \frac{\text{false wakes allowed}}{\text{windows listened}}}%
  {{\raggedright One wake a month over one-second windows: $1/\val{fwwin} \approx \val{fwfpr}$.
   With zero false wakes observed, hours needed $\approx \val{fwthree} \div$ allowed
   false wakes per hour (\val{fwconf}\% confidence), so one per \val{fwmonthh}-hour month needs
   \val{fwhours} hours of representative non-wake audio.\par}}{}{}

\vspace{11pt}
\sheetsec{The pipeline and its two bounds}
{\centering
\begin{tikzpicture}[sheet, y=1.25pt, stg/.style={accbox, minimum width=40pt, minimum height=17pt}]
  \foreach \i/\t in {0/Acquire,1/Validate,2/Label,3/Process,4/Release}{
    \node[stg] (c\i) at ({24+\i*50.5},62) {\t};}
  \foreach \i in {0,...,3}{\pgfmathtruncatemacro\j{\i+1}\draw[arr] (c\i) -- (c\j);}
  \node[note, anchor=south] at (125,73) {capability: what the model can learn};
  \node[nbox, minimum width=60pt, minimum height=17pt] (s0) at (40,14) {Storage};
  \node[nbox, minimum width=60pt, minimum height=17pt] (s1) at (125,14) {Host decode};
  \node[hotbox, minimum width=60pt, minimum height=17pt] (s2) at (210,14) {Accelerator};
  \draw[arr] (s0) -- (s1); \draw[arr] (s1) -- (s2);
  \draw[arr, rounded corners=3pt] (c4.south) -- (226,40) -- (40,40) -- (s0.north);
  \node[note, anchor=north] at (125,3) {throughput: how fast it learns};
\end{tikzpicture}\par}
\end{rightcol}

% ===================================================================== BACK: PRACTICE
\backside
\begin{answers}
\ans{Finish}{\val{disk} MB/s $\div$ \val{deckb} KB $=$ \val{ygot} img/s; $\val{one} - \val{ygot}/\val{rate} =$ \val{yloss}\%
  lost, worse than \val{loss}\%. The cache removed decoding but multiplied bytes per sample by \val{expand},
  and fixed the stored crops. (Supply bound)}
\ans{1}{$\val{one} \div \val{qwin} = \val{qfpr}$ per window. Hours needed: $\val{fwthree} \div (\val{one}$ per
  \val{qweekh} h$) =$ \val{qhours} h, \val{qweeks} weeks with no false wake. (False-wake budget)}
\ans{2}{Per-clip peak uses the finished clip, which a streaming device never has: skew. Whole-corpus
  fitting leaks test statistics into training, so the test score is no longer trustworthy. Fit on the training split, store,
  replay. (Fit once, replay everywhere)}
\columnbreak
\ans{3}{(a) Decode-bound: Snappy, less work per byte. (b) Storage-bound: gzip, fewer bytes per sample
  raise the supply rate. (Supply bound)}
\ans{4}{\val{psiobs} $>$ \val{psicrit}: new-microphone inputs moved; it says nothing yet about accuracy.
  Replay fixed audio through old and new preprocessing; rescore fixed predictions under both rubrics;
  compare microphone cohorts, matched for speakers and rooms, on one path and one rubric. (Which failure)}
\end{answers}

\begin{leftcol}
\sheetsec{Worked example}
\begin{practicebox}\boxtitle{Can a cloud disk keep one accelerator fed?}
ResNet-50 training sustains \val{rate} images/s on one accelerator. Each image is a \val{jpgkb} KB JPEG
read from disk every epoch; the volume supplies \val{disk} MB/s.\par\vspace{5pt}
{\centering
\begin{tikzpicture}[sheet]
  \pgfmathsetmacro\u{150/\val{need}}
  \node[note, anchor=east] at (52,30) {needed};
  \draw[nbox] (56,23) rectangle ({56+\val{need}*\u},37);
  \node[lab, anchor=west] at ({60+\val{need}*\u},30) {\val{need} MB/s};
  \node[note, anchor=east] at (52,8) {supplied};
  \draw[hotbox] (56,1) rectangle ({56+\val{disk}*\u},15);
  \node[lab, anchor=west] at ({60+\val{disk}*\u},8) {\val{disk} MB/s};
\end{tikzpicture}\par}\vspace{4pt}
\begin{steps}
\item Demand: \val{rate} img/s $\times$ \val{jpgkb} KB $=$ \res{\val{need} MB/s}.
\item Supply: \val{disk} MB/s $\div$ \val{jpgkb} KB $=$ \res{\val{got} img/s}, the smaller rate, so it sets the pace.
\item Feeding loss: $\val{one} - \val{got}/\val{rate} =$ \res{\val{loss}\%} of the sustained rate.
\end{steps}
\takeaway{Compare bytes needed per second with bytes supplied; the smaller rate wins, however fast the accelerator.}
\end{practicebox}

\vspace{12pt}
\sheetsec{Your turn: finish it}
\begin{practicebox}\boxtitle{Cache decoded tensors to skip decoding?}
To spare the host, the team stores each image already decoded, as a 32-bit tensor of \val{deckb} KB
(\val{expand}$\times$ the JPEG), on the same \val{disk} MB/s volume.\par\vspace{3pt}
\textbf{Given:} demand $=$ \val{rate} $\times$ \val{deckb} KB $=$ \val{yneed} GB/s.\par
\textbf{You:} find the supplied rate and the feeding loss, and say what the cache traded away.
\end{practicebox}
\end{leftcol}%
\begin{rightcol}
\sheetsec{Common mistakes}
\mistake{``More data always improves the model.''}%
  {In the keyword-spotting design, \val{morex}$\times$ the examples buys \val{moregainlo}--\val{moregainhi}\%
   more detection for \val{morex}$\times$ the labeling cost. Redundant samples still cost I/O.}{}
\mistake{``Training and serving share the feature code, so they match.''}%
  {If serving recomputes the mean on each batch, the same code feeds the model different numbers.
   Store the fitted state with the model.}{}
\mistake{``Our logs show where the wake-word model fails.''}%
  {The device uploads audio only when it wakes, so a missed wake word leaves no recording.
   Measure misses with separate, consented collection.}{}

\sheetsec{Check yourself}
\begin{qs}
\item \assumed{} a budget of one false wake per week, one-second windows, always listening. What
  per-window false-positive rate is needed? If none occur, how many hours of representative
  non-wake audio support the claim?
\item The detector decides as audio streams in, before a clip ends. One engineer scales each stored clip to its own peak loudness; another fits mean and variance on
  the whole corpus before splitting. What goes wrong in each case, and what is the fix?
\item A loader starves the accelerator. (a) Host CPUs at full load, disk half busy. (b) Disk at full
  bandwidth, CPUs idle. Snappy (decodes fast, compresses less) or gzip (compresses more, decodes
  slowly)? Why?
\item \val{q4months} months after launch, a wake-word detector's measured detection rate is \val{q4drop}
  points lower. A new microphone shipped, the grading rubric was revised, and preprocessing changed.
  \assumed{} PSI of new-microphone inputs vs.\ training is \val{psiobs}. What does that tell you, and what three
  comparisons separate the causes?
\end{qs}
\end{rightcol}
\end{document}
